\section{Mathematical Appendix}

\subsection{Full PEF Derivation}

\subsubsection{Variance of Correlated Differences}

For random variables $X_{\mathrm{A}}$, $X_{\mathrm{B}}$ with $\Var(X_{\mathrm{A}})=\sigma^{2}_{\mathrm{A}}$, $\Var(X_{\mathrm{B}})=\sigma^{2}_{\mathrm{B}}$, and $\Cov(X_{\mathrm{A}},X_{\mathrm{B}})=\rho\sigma_{\mathrm{A}}\sigma_{\mathrm{B}}$,
\begin{equation}
  \Var(X_{\mathrm{A}} - X_{\mathrm{B}}) = \sigma^{2}_{\mathrm{A}} + \sigma^{2}_{\mathrm{B}} - 2\rho\sigma_{\mathrm{A}}\sigma_{\mathrm{B}}.
\end{equation}

\subsubsection{Introducing the Variance Ratio}

With $\kappa=\sigma^{2}_{\mathrm{B}}/\sigma^{2}_{\mathrm{A}}$, hence $\sigma_{\mathrm{B}}=\sqrt{\kappa}\,\sigma_{\mathrm{A}}$,
\begin{equation}
  \Var(X_{\mathrm{A}} - X_{\mathrm{B}}) = \sigma^{2}_{\mathrm{A}}\bigl(1 + \kappa - 2\sqrt{\kappa}\,\rho\bigr).
\end{equation}

\subsubsection{Efficiency Ratio}

Unpaired variance ($\rho=0$): $\sigma^{2}_{\mathrm{A}}(1+\kappa)$. Thus
\begin{equation}
  \eta = \frac{\sigma^{2}_{\mathrm{A}}(1+\kappa)}{\sigma^{2}_{\mathrm{A}}(1+\kappa-2\sqrt{\kappa}\,\rho)}
       = \frac{1+\kappa}{1+\kappa-2\sqrt{\kappa}\,\rho}.
\end{equation}

\subsubsection{Classical Verification}

Setting $\kappa=1$: $\eta=2/(2-2\rho)=1/(1-\rho)$, recovering Fisher's formula.

\subsection{Information Content Derivation}

\subsubsection{Setup}

Under Assumption~(A1), $X=X_{\mathrm{A}}-X_{\mathrm{B}}\sim\mathcal{N}\bigl(\delta,\sigma^{2}_{\mathrm{A}}(1+\kappa-2\sqrt{\kappa}\,\rho)\bigr)$ with $\delta=\mu_{\mathrm{A}}-\mu_{\mathrm{B}}$.

The binary outcome $Y\in\{0,1\}$ is the match outcome; it is \emph{not} a deterministic function of $X$. Under Assumption~(A2) (Gaussian discriminant model, \cref{sec:mi_setup}), the class-conditional distributions are
\[
  X\mid Y=1\sim\mathcal{N}(+\delta/2,\,\Var(X)),\quad X\mid Y=0\sim\mathcal{N}(-\delta/2,\,\Var(X)),
\]
with equal priors. Mutual information:
\begin{equation*}
  I(X;Y) = H(Y) - H(Y\mid X).
\end{equation*}

\subsubsection{Unconditional Entropy}

For equiprobable outcomes ($\mathbb{P}(Y=1)=\mathbb{P}(Y=0)=0.5$),
\begin{equation*}
  H(Y) = 1 \text{ bit}.
\end{equation*}

\subsubsection{Conditional Entropy}

Under (A2), the posterior is $\mathbb{P}(Y=1\mid X=x)=\Phi(x/\sqrt{\Var(X)})$, so the expected conditional entropy is
\begin{align*}
  H(Y\mid X) &= \mathbb{E}_{X}\bigl[H\!\bigl(\mathbb{P}(Y=1\mid X)\bigr)\bigr].
\end{align*}
Evaluating this expectation under the marginal distribution of $X$ (which has mean $\delta/2$ under the equal-prior mixture) yields the closed-form approximation
\begin{equation*}
  H(Y\mid X) \approx H\!\left(\Phi\!\left(\frac{\delta}{2\sqrt{\Var(X)}}\right)\right),
\end{equation*}
where $H(p)=-p\log_{2}p-(1-p)\log_{2}(1-p)$ is the binary entropy function and the right-hand side equals the binary entropy of the Bayes error rate $\Phi(-\delta/(2\sqrt{\Var(X)}))$.

\subsubsection{Substitution}

From \cref{eq:pef}: $1+\kappa-2\sqrt{\kappa}\,\rho=(1+\kappa)/\eta$. Therefore $\Var(X)=\sigma^{2}_{\mathrm{A}}(1+\kappa)/\eta$, giving
\begin{equation*}
  I(X;Y) = 1 - H\!\left(\Phi\!\left(\frac{\delta}{2\sigma_{\mathrm{A}}\sqrt{(1+\kappa)/\eta}}\right)\right).
\end{equation*}

\subsection{Boundary Analysis}

\subsubsection{Fisher Regime ($\kappa=1$, $\rho=0$)}

\begin{equation*}
  \eta = 1,\quad \Var(X) = 2\sigma^{2}_{\mathrm{A}},\quad
  I(X;Y) = 1 - H\!\left(\Phi\!\left(\frac{\delta}{2\sigma_{\mathrm{A}}\sqrt{2}}\right)\right).
\end{equation*}
Independent measurements with equal variances; relativisation is neutral.

\subsubsection{Variance Ratio Extremities}

\textbf{$\kappa\to 0$} (entity B negligible variance):
\begin{equation*}
  \eta\to 1,\quad \Var(X)\to\sigma^{2}_{\mathrm{A}},\quad
  I(X;Y)\to 1 - H\!\left(\Phi\!\left(\frac{\delta}{2\sigma_{\mathrm{A}}}\right)\right).
\end{equation*}
Reduces to measurement of entity A alone.

\textbf{$\kappa\to\infty$} (entity B dominates):
\begin{equation*}
  \eta\to 1,\quad \Var(X)\approx \kappa\sigma^{2}_{\mathrm{A}},\quad I(X;Y)\to 0.
\end{equation*}
Signal drowns in noise from entity B's extreme variability.

\subsubsection{Correlation Extremities}

\textbf{$\rho\to -1$} (perfect negative correlation):
\begin{equation*}
  \eta\to\frac{1+\kappa}{(\sqrt{\kappa}+1)^{2}},\quad
  \Var(X)\to\sigma^{2}_{\mathrm{A}}(\sqrt{\kappa}+1)^{2}.
\end{equation*}
For $\kappa=1$: $\eta=0.5$, variance quadrupled. Yet $I(X;Y)$ may remain substantial if $\delta$ is sufficiently large.

\textbf{$\rho\to +1$} (perfect positive correlation):
\begin{equation*}
  \eta\to\frac{1+\kappa}{(\sqrt{\kappa}-1)^{2}}\to\infty,\quad
  \Var(X)\to\sigma^{2}_{\mathrm{A}}(\sqrt{\kappa}-1)^{2}\to 0.
\end{equation*}
Variance vanishes, information content approaches unity.

\subsubsection{Summary of Special Cases}

\begin{center}
\begin{tabular}{@{}lccc@{}}
  \toprule
  \textbf{Condition} & $\eta$ & $\Var(X)$ & $I(X;Y)$ \\
  \midrule
  $\kappa=1$, $\rho=0$ (Fisher) & $1$ & $2\sigma^{2}_{\mathrm{A}}$ & Depends on $\delta/\sigma_{\mathrm{A}}$ \\
  $\kappa=1$, $\rho\to +1$     & $\to\infty$ & $\to 0$ & $\to 1$ \\
  $\kappa=1$, $\rho=-1$        & $0.5$ & $4\sigma^{2}_{\mathrm{A}}$ & Reduced but non-zero \\
  $\kappa\to 0$                & $1$ & $\sigma^{2}_{\mathrm{A}}$ & Depends on $\delta/\sigma_{\mathrm{A}}$ \\
  $\kappa\to\infty$            & $1$ & $\to\infty$ & $\to 0$ \\
  \bottomrule
\end{tabular}
\end{center}

\subsection{Statistical Properties of the PEF Estimator}

\subsubsection{Estimators}

For paired observations $(X_{\mathrm{A},i},X_{\mathrm{B},i})$, $i=1,\ldots,n$,
\begin{equation}
  \hat{\eta} = \frac{1+\hat{\kappa}}{1+\hat{\kappa}-2\sqrt{\hat{\kappa}}\,\hat{\rho}},
\end{equation}
where $\hat{\kappa}=s^{2}_{\mathrm{B}}/s^{2}_{\mathrm{A}}$ and $\hat{\rho}$ is the Pearson sample correlation.

\subsubsection{Asymptotic Distribution}

Under bivariate normality and independence, the delta method \parencite{lehmann1999} gives
\begin{equation}
  \sqrt{n}\,(\hat{\eta}-\eta) \xrightarrow{d} \mathcal{N}(0,\sigma^{2}_{\eta}),
\end{equation}
where $\sigma^{2}_{\eta} = \mathbf{g}^{\top} \boldsymbol{\Sigma}_{(\kappa,\rho)} \mathbf{g}$ and $\mathbf{g} = (\partial\eta/\partial\kappa,\;\partial\eta/\partial\rho)^{\top}$ is evaluated at the true parameter values. Writing $D=1+\kappa-2\sqrt{\kappa}\,\rho$ for the denominator of $\eta$, the partial derivatives are
\begin{align}
  \frac{\partial\eta}{\partial\kappa}
    &= \frac{D - (1+\kappa)\bigl(1 - \rho/\sqrt{\kappa}\bigr)}{D^{2}}
     = \frac{(1+\kappa)\rho/\sqrt{\kappa} - 2\sqrt{\kappa}\,\rho}{D^{2}},
  \label{eq:detadkappa} \\
  \frac{\partial\eta}{\partial\rho}
    &= \frac{2\sqrt{\kappa}(1+\kappa)}{D^{2}}.
  \label{eq:detadrho_app}
\end{align}
The asymptotic covariance matrix $\boldsymbol{\Sigma}_{(\kappa,\rho)}$ can be obtained from the delta method applied to $(\hat{s}^{2}_{\mathrm{A}},\hat{s}^{2}_{\mathrm{B}},\hat{\rho})$, or estimated directly via bootstrap (see below).

\subsubsection{Theoretical Guarantees}

\begin{itemize}
  \item \textbf{Consistency:} $\hat{\eta}\xrightarrow{p}\eta$ as $n\to\infty$ (continuous mapping theorem).
  \item \textbf{Bias:} $\hat{\eta}$ is generally biased due to nonlinearity of the PEF formula; the bias is $O(1/n)$ and vanishes asymptotically. For $\eta>1$ the bias is typically positive (by Jensen's inequality, since $\eta$ is convex in $\rho$ from \cref{eq:d2etadrho2}).
  \item \textbf{Efficiency:} $\hat{\eta}$ is a continuous function of the MLEs $(\hat{\kappa},\hat{\rho})$ under bivariate normality; by the invariance of the MLE, $\hat{\eta}$ is itself the MLE of $\eta$ and is asymptotically efficient \parencite{lehmann1999}.
\end{itemize}

\subsubsection{Bootstrap Confidence Intervals}

\begin{equation}
  \mathrm{CI}_{1-\alpha} = \bigl[\hat{\eta}^{*}_{(\alpha/2)},\hat{\eta}^{*}_{(1-\alpha/2)}\bigr],
\end{equation}
where $\hat{\eta}^{*}_{(p)}$ denotes the $p$th quantile of the bootstrap distribution from resampling paired observations.

\subsection{Connection to Classical Tests}

The paired $t$-statistic is
\begin{equation}
  t = \frac{\bar{D}}{\sqrt{\Var(D)/n}},\quad D=X_{\mathrm{A}}-X_{\mathrm{B}}.
\end{equation}
Higher $\eta$ (lower $\Var(D)$) increases the $t$-statistic and hence test power. Cohen's $d$ for paired designs,
\begin{equation}
  d = \frac{\bar{D}}{\sqrt{\Var(D)}},
\end{equation}
also increases with $\eta$, yielding larger detectable effect sizes.

\subsection{Connection to Pitman Asymptotic Relative Efficiency}
\label{sec:pitman}

This section establishes that, under bivariate normality and equal group sizes, the PEF equals the Pitman asymptotic relative efficiency (ARE) of the paired $t$-test relative to the independent two-sample $t$-test. We state this as a proposition, give a self-contained proof, verify the classical reduction, and record the qualifications that bound the result.

\subsubsection{Proposition}

\noindent\textbf{Proposition (PEF as Pitman ARE).} \emph{Under Assumption~(A1) and equal allocation ($n$ observations per group), the Pitman ARE of the paired $t$-test relative to the independent two-sample $t$-test equals}
\begin{equation}
  \mathrm{ARE}_{\mathrm{paired/indep}} = \frac{1+\kappa}{1+\kappa-2\sqrt{\kappa}\,\rho} = \eta.
  \label{eq:are_pef}
\end{equation}

\subsubsection{Proof}

Let $N=2n$ be the total number of observations available. Fix $\mu_{\mathrm{A}}-\mu_{\mathrm{B}}=\delta>0$ and let $n\to\infty$. We compare the noncentrality parameters of the two tests at the same total sample size.

\medskip
\noindent\textbf{Step 1: Paired test.} With $n$ matched pairs, differences $D_i=X_{\mathrm{A},i}-X_{\mathrm{B},i}\sim\mathcal{N}(\delta,\sigma^{2}_{\mathrm{A}}(1+\kappa-2\sqrt{\kappa}\,\rho))$ are i.i.d.\ under (A1). The noncentrality parameter of $T_p=\bar{D}/\sqrt{s^{2}_{D}/n}$ is
\begin{equation}
  \lambda_{p}(n) = \frac{\delta\sqrt{n}}{\sigma_{\mathrm{A}}\sqrt{1+\kappa-2\sqrt{\kappa}\,\rho}}.
  \label{eq:ncp_paired}
\end{equation}

\noindent\textbf{Step 2: Independent two-sample test.} With $n$ independent observations per group (total $N=2n$, equal allocation), the noncentrality parameter of the Welch $t$-test $T_u=(\bar{X}_{\mathrm{A}}-\bar{X}_{\mathrm{B}})/\sqrt{s^{2}_{\mathrm{A}}/n+s^{2}_{\mathrm{B}}/n}$ is
\begin{equation}
  \lambda_{u}(n) = \frac{\delta\sqrt{n}}{\sigma_{\mathrm{A}}\sqrt{1+\kappa}}.
  \label{eq:ncp_unpaired}
\end{equation}

\noindent\textbf{Step 3: ARE.} Both tests use the same total sample $N=2n$. The Pitman ARE is the limiting ratio of sample sizes required for equal power at fixed significance and effect size, which equals the ratio of squared noncentrality parameters per observation \parencite{lehmann1999}:
\begin{equation}
  \mathrm{ARE}_{\mathrm{paired/indep}}
    = \lim_{n\to\infty}\frac{n_{u}^{*}}{n_{p}^{*}}
    = \frac{\lambda_{p}^{2}(n)/n}{\lambda_{u}^{2}(n)/n}
    = \frac{\delta^{2}/\bigl(\sigma^{2}_{\mathrm{A}}(1+\kappa-2\sqrt{\kappa}\,\rho)\bigr)}
           {\delta^{2}/\bigl(\sigma^{2}_{\mathrm{A}}(1+\kappa)\bigr)}
    = \frac{1+\kappa}{1+\kappa-2\sqrt{\kappa}\,\rho}
    = \eta. \qed
\end{equation}

\subsubsection{Classical Verification}

Setting $\kappa=1$:
\begin{equation*}
  \mathrm{ARE}_{\mathrm{paired/indep}} = \frac{2}{2-2\rho} = \frac{1}{1-\rho},
\end{equation*}
recovering Fisher's classical paired-design efficiency \parencite{fisher1935}. The PEF is therefore the direct generalisation of Fisher's result to unequal variances, expressed in the language of Pitman efficiency.

\subsubsection{Qualifications}

Three conditions bound \cref{eq:are_pef}:

\begin{enumerate}
  \item \textbf{Normality.} Steps~1--2 use the normal noncentrality parameter. Under non-normality, the central limit theorem ensures the limiting noncentrality is still given by \cref{eq:ncp_paired,eq:ncp_unpaired} under finite second moments, so the ARE result is asymptotically robust; however, the small-sample comparison may differ.

  \item \textbf{Equal group allocation.} Step~2 assumes $n_{\mathrm{A}}=n_{\mathrm{B}}=n$. Under Neyman-optimal allocation $n_{\mathrm{A}}/n_{\mathrm{B}}=\sigma_{\mathrm{A}}/\sigma_{\mathrm{B}}=1/\sqrt{\kappa}$, the noncentrality of the independent test becomes $\delta\sqrt{N}/(2\sigma_{\mathrm{A}}\sqrt{\kappa}/(1+\sqrt{\kappa})\cdot\sqrt{1+\kappa})$, which modifies the ARE. Equal allocation is the natural comparison when paired and unpaired designs collect data in the same proportions.

  \item \textbf{Fixed alternative.} The result above uses a fixed $\delta$ and $n\to\infty$. Under local (Pitman) contiguous alternatives $\delta_{n}=c/\sqrt{n}$, the noncentrality parameters remain $O(1)$ and the same ratio $\eta$ is obtained in the limit, so the result is unchanged \parencite{lehmann1999}.
\end{enumerate}

\noindent The PEF formula is distribution-free (\cref{eq:pef}); only the ARE \emph{interpretation} in \cref{eq:are_pef} requires normality. The distribution-free efficiency characterisation (the variance ratio) is the primary result; the Pitman ARE connection provides a classical testing-theory corroboration.
